\chapter{Representation of Physical Quantities on the Semantic Web}

This part examines the representation problem introduced in Section~1. 
Section~3.1 presents the ontology-based approaches that constitute current practice. 
Section~3.2 analyses their limitations in detail, ordered from the most fundamental to the least. 
Section~3.3 introduces \ac{ucum}, the code system whose grammar resolves the coverage problem structurally. 
Section~3.4 presents the \ac{cdt} specification, which builds on \ac{ucum} to address the  problems at the datatype level.

\section{Ontology-Based Approaches}

The standard approach to representing physical quantities on the Semantic Web is ontology-based. 
Several ontologies have been developed for this purpose. 
Keil and Schindler evaluated eight of them: \ac{muo}, \ac{oboe}, \ac{om}, \ac{qu}, \ac{qudt}, \ac{sweet}, \ac{uo}, and Wikidata \cite{keil2018comparison}. 
\ac{qudt} and \ac{om}, examined in detail below, share one structural commitment: a single quantity value requires either a named individual or a blank node, together with several triples.

\ac{qudt}, originating from a \ac{nasa} initiative, organises units within a hierarchy of quantity kinds and dimensions \cite{hodgson2014quantities}. 
In \ac{qudt}, a quantity value is represented as a resource typed as \textsf{qudt:QuantityValue}. 
The magnitude is carried by \textsf{qudt:numericValue} as a numeric \ac{xsd} literal, such as \textsf{xsd:double}, \textsf{xsd:decimal}, or \textsf{xsd:integer}; the unit is referenced by \textsf{qudt:unit} as a named individual. 
The following Turtle fragment illustrates the current \ac{qudt} encoding pattern:

\begin{verbatim}
@prefix sosa:  <http://www.w3.org/ns/sosa/> .
@prefix xsd:   <http://www.w3.org/2001/XMLSchema#> .
@prefix rdfs:  <http://www.w3.org/2000/01/rdf-schema#> .
@prefix qudt:  <http://qudt.org/schema/qudt/> .
@prefix unit:  <http://qudt.org/vocab/unit/> .

<Observation/234534> a sosa:Observation ;
   rdfs:comment "Observation of the difference between the
        outside temperature and the inside temperature."@en ;
   sosa:hasFeatureOfInterest <apartment/134> ;
   sosa:hasResult [
      a qudt:QuantityValue ;
      qudt:unit unit:DEG_C ;
      qudt:numericValue "-29.9"^^xsd:double ] .
\end{verbatim}

Four triples encode one quantity value. 
The magnitude and the unit are separated across two independent statements on the blank node, and no shared type binds them together. \ac{qudt} also encodes dimensional structure: each unit carries a \textsf{qudt:hasDimensionVector} property linking it to a vector of exponents over the fundamental dimensions, and a \textsf{qudt:conversionMultiplier} relating it to its \ac{si} base equivalent \cite{hodgson2014quantities}.

\ac{om} , described by Rijgersberg, van Assem, and Top, follows a comparable reification pattern with finer-grained unit groupings \cite{vanassem2013ontology}. 
A quantity value is an \textsf{om:Measure} blank node carrying a numeric value and a unit reference. 
An equivalent encoding in \ac{om} is:

\begin{verbatim}
@prefix sosa: <http://www.w3.org/ns/sosa/> .
@prefix xsd:  <http://www.w3.org/2001/XMLSchema#> .
@prefix rdfs: <http://www.w3.org/2000/01/rdf-schema#> .
@prefix om:   <http://www.ontology-of-units-of-measure.org/resource/om-2/> .

<Observation/234534> a sosa:Observation ;
   rdfs:comment "Observation of the difference between the 
        outside temperature and the inside temperature."@en ;
   sosa:hasFeatureOfInterest <apartment/134> ;
   sosa:hasResult [
      a om:Measure ;
      om:hasUnit om:degreeCelsius ;
      om:hasNumericalValue "-29.9"^^xsd:double ] .
\end{verbatim}

The structural commitment is the same: a blank node, a type triple, a magnitude triple, and a unit triple. 
\ac{om} provides more refined coverage than \ac{qudt} and finer-grained quantity kind hierarchies, but the reification pattern does not change. 
\ac{om} also encodes dimensional structure, through properties such as \textsf{SI\_mass\_dimension\_exponent} and \textsf{SI\_length\_dimension\_exponent} on dimension individuals \cite{vanassem2013ontology}.

SOSA and SSN, published as a joint W3C and OGC Recommendation, take an observation-centric approach \cite{haller2019sosa}. A sensor observation is modelled as a sosa:Observation linked to a result via sosa:hasResult. SOSA and SSN do not prescribe how the numeric value and unit are encoded within the result; they are typically combined with QUDT or OM for that purpose. The observation structure therefore adds further triples ( as can be seen in the previous examples) on top of the quantity value pattern from whichever unit ontology is used.

Several further ontologies address the same problem space. 
\ac{muo} originated as a project to apply semantics in mobile environments, with its individuals automatically generated from \ac{ucum} \cite{keil2018comparison}. 
\ac{oboe} takes an observation-centric approach developed primarily for ecological and scientific measurement scenarios \cite{madin2007ontology}. 
\ac{qu} is a showcase ontology based on the \ac{omg} SysML 1.2 \ac{qudv} specification and the \ac{uncefact} Recommendation 20 code list, published through the \ac{w3c} Semantic Sensor Network Incubator Group \cite{keil2018comparison}. 
\ac{sweet} covers a broad earth science vocabulary that includes units of measurement \cite{raskin2005knowledge}. 
\ac{uo}, part of the \ac{obo} Foundry, models units alongside the \ac{pato} \cite{keil2018comparison}.

\section{Limitations of Ontology-Based Approaches}

The most fundamental problem is computational. 
\ac{qudt} represents dimensional structure through a vector of exponents over the fundamental dimensions: length, mass, time, electric current, temperature, amount of substance, and luminous intensity \cite{hodgson2014quantities}. 
It provides \textsf{qudt:conversionMultiplier} per unit, relating each unit to its \ac{si} base equivalent. 
\ac{om} models dimensions similarly, with properties such as \textsf{SI\_mass\_dimension\_exponent} and \textsf{SI\_length\_dimension\_exponent} on dimension individuals \cite{vanassem2013ontology}. 
This information is present in the graph and queryable in principle. 
It does not reach the \ac{sparql} expression evaluator.

A FILTER expression evaluates against the \ac{rdf} terms bound to query variables in the current solution, as established in Section~2.3. 
For a \ac{qudt}-encoded quantity, the term bound to the magnitude variable is a bare \textsf{xsd:double}, while the unit \ac{iri} is bound separately through a separate triple pattern. 
They are not a composite typed value. 
The dimensional relationships in the ontology are not carried into expression evaluation alongside the magnitude. 
A FILTER compares the raw \textsf{xsd:double} magnitude against 1.0 km with no awareness of units. 
A measurement of 500 metres passes FILTER(500 > 1.0) and is incorrectly included in the results, despite being shorter than one kilometre. 
Lefrançois and Zimmermann confirm that no existing \ac{rdf} or \ac{sparql} engine 
provides unit-aware evaluation for \ac{qudt} or \ac{om} quantities \cite{lefranccois2018unified}. 
At the time of that work, no such engine existed, and to the best of our knowledge 
that remains the case today.
Workarounds require making conversion factors explicit in every unit-aware query. 
When conversion information is retrieved from the ontology rather than hardcoded, additional triple patterns joining across the unit graph are needed, compounding the complexity of every unit-aware query.

The second problem is bounded coverage. 
Ontology-based unit vocabularies enumerate units as named individuals in a curated graph, and the scale of the enumeration required is larger than it appears. 
Length alone spans metric, customary, nautical, and astronomical systems. 
Derived units expand the space further: every pairing of a length unit with a time unit is a distinct speed unit, and the same holds for pressure, energy, and power. 
Each additional physical dimension multiplies the number of valid unit expressions by the number of units in that dimension. 
The set of valid unit expressions is unbounded.

Specific gaps have been documented across the evaluated ontologies. 
Keil and Schindler, evaluating eight major unit ontologies, found that 
coverage differed substantially between systems \cite{keil2018comparison}. 
When a needed unit is absent, the extension mechanism requires defining 
it by its base units and conversion factors. 
Two teams independently extending the vocabulary for the same missing 
unit create two different \acp{iri} for the same physical unit. 
The coverage gap thereby becomes an interoperability gap \cite{lefranccois2018unified}.

The third problem is interoperability. 
Multiple independently developed unit ontologies create a structural barrier at the \ac{iri} level. 
Keil and Schindler evaluated eight ontologies: \ac{muo}, \ac{oboe}, \ac{om} 1, \ac{om} 2, \ac{qu}, \ac{qudt}, \ac{sweet}, and \ac{uo}, together with the Wikidata corpus \cite{keil2018comparison}. 
They found that these systems use different terms for the same concepts. 
The concept \ac{vim} calls kind of quantity appears as physical quality in \ac{muo}, quantity kind in \ac{qu} and \ac{qudt}, and physical quantity in Wikidata. 
Unit individuals receive different \acp{iri} across all eight systems, and mapping them revealed sparse \textsf{owl:sameAs} links and incorrect conversion factors in several ontologies \cite{keil2018comparison}. 
Datasets encoded against different ontologies, or against different versions of the same ontology, cannot be queried uniformly without cross-ontology alignments. 
Those alignments are manually constructed and are not guaranteed to be present.

The fourth problem is verbosity. 
The reification pattern shown in Section~3.1 requires at minimum four 
triples per quantity value regardless of content: one linking the subject 
to a blank node or named individual, one typing that resource, one 
carrying the magnitude, and one carrying the unit reference. 
Lefrançois and Zimmermann note that complex canonicalization mechanisms 
are needed to query these values uniformly \cite{lefranccois2018unified}. 
In a sensor network generating thousands of observations, this overhead 
compounds proportionally with the dataset size. 


\section{UCUM}

The \ac{ucum} is a code system developed by Gunther Schadow and Clem McDonald and distributed by the Regenstrief Institute \cite{schadow_ucum}. 
Its purpose is to enable unambiguous machine-to-machine communication of quantities together with their units. 
Natural language unit names vary across languages, communities, and domains. 
\ac{ucum} expressions, by contrast, are case-sensitive character strings with a formal grammar, which makes them suitable for electronic data interchange without a separate disambiguation layer.

The grammar is compact and precisely defined. 
A unit expression is built from atomic unit symbols and metric prefixes using two operators, a period (.) for multiplication and a solidus (/) for division, with signed integer exponents written directly after the symbol. 
All symbols are case-sensitive: m denotes the metre and k the kilo prefix. 
Non-metric legacy units are enclosed in square brackets, such as \textsf{[in\_i]} for the international inch. 
Dimensionless annotations are enclosed in curly braces. 
The grammar assigns each valid \ac{ucum} expression a precisely defined dimensional interpretation.

Coverage spans scientific, engineering, and business domains. 
\ac{si} base and derived units are included alongside metric prefixes, \ac{us} customary units, and domain-specific units used in laboratory reporting, radiology, and pharmacology. 
\ac{ucum} has been adopted by \ac{dicom} \cite{nema_dicom}, \ac{loinc} \cite{mcdonald2003loinc}, and \ac{hl7} \cite{hl7_fhir_quantity}, and is referenced by \ac{iso} 11240:2012 \cite{iso2012_11240}.

\ac{ucum} defines a subset of its units as special units, whose conversion to base-unit form requires a non-linear function. 
Celsius (Cel) and Fahrenheit (\textsf{[degF]}) involve an additive offset relative to the thermodynamic temperature scale. 
pH involves a logarithmic transformation of hydrogen ion concentration. 
The bel (B) and its decimal subdivision the decibel are logarithmic ratios. 
These units are valid \ac{ucum} expressions, but their conversion semantics fall outside the multiplicative algebra.

\section{The CDT Specification}

Lefrançois and Zimmermann presented \textsf{cdt:ucum} at ESWC 2018 and formalized its
specification in the LINDT Custom Datatypes document \cite{lefranccois2018unified,lefrancois2021lindt}.
A quantity value is encoded as a single typed literal whose unit is part of the lexical form
rather than a separate triple. Operator semantics are defined for equality, comparison, and
arithmetic, making them available through native SPARQL operators without custom function syntax.

The specification defines two general-purpose datatypes: \textsf{cdt:ucum}, with IRI
\texttt{https://\allowbreak w3id.org/\allowbreak cdt/\allowbreak ucum}, and
\textsf{cdt:ucumunit}, with IRI
\texttt{https://\allowbreak w3id.org/\allowbreak cdt/\allowbreak ucumunit}. A quantity value
encoded with \textsf{cdt:ucum} requires a single triple:

\begin{verbatim}
@prefix cdt: <https://w3id.org/cdt/> .
BASE <http://example.org/>

<observation/235727> a sosa:Observation ;
    sosa:hasFeatureOfInterest <house/134/kitchen> ;
    sosa:observedProperty <electricConsumption> ;
    sosa:madeBySensor <sensor/927> ;
    sosa:hasSimpleResult "22.4 kWh"^^cdt:ucum .

\end{verbatim}

Compared to the four-triple QUDT pattern in Section~3.1, the magnitude and unit are part of
a single typed literal and no blank node is needed.

\textsf{cdt:ucumunit} encodes a bare UCUM unit expression without a magnitude. Its lexical
space is the set of valid UCUM unit expressions. Its value space is the underlying set of the
algebra of units defined in Section~3 of the UCUM specification. That algebra is formalized
in §18, which defines the set $U$ of units as the equivalence classes generated by the
equality relation over UCUM expressions; two expressions belong to the same equivalence class
when they denote the same unit, and $(U, \cdot)$ forms an Abelian group under unit
multiplication. The lexical-to-value mapping sends each unit expression to its equivalence
class in $U$. Two \textsf{cdt:ucumunit} literals are value-equal when they denote the same
unit: \texttt{km.h-1}, \texttt{km/h}, and \texttt{(1000m)/(60min)} are all well-formed and
share the same value.

\textsf{cdt:ucum} encodes a quantity value as a single literal. Its lexical space is the
concatenation of the lexical form of an \textsf{xsd:decimal}, optionally followed by
\texttt{e} or \texttt{E} and the lexical form of an \textsf{xsd:integer}, at least one
whitespace character, and a valid UCUM unit expression. The value space is the set of pairs
$(v, u)$ where $v$ is a real number and $u$ is a product of UCUM base units raised to integer
powers. This follows from §20 of the UCUM specification, which represents every proper unit
as a pair of a scalar magnitude and a dimension vector over the base units; in the value pair,
the magnitude is absorbed into $v$ and $u$ retains only the base-unit structure. The
lexical-to-value mapping normalizes both components: the number $v = m \times 10^e$, where
$m$ is the decimal and $e$ is the exponent (zero if absent); the unit $u = r \cdot b$, where
$r$ is the magnitude of the expressed unit relative to the base units $b$. The mapping assigns
the lexical form to the pair $(v \times r,\, b)$.

Value equality follows directly from the mapping. The literal
\texttt{"1 km"\^{}\^{}}\textsf{cdt:ucum} has decimal~1 and the unit \texttt{km}, whose
magnitude relative to the metre is~1000. Its mapped value is
$(1 \times 1000,\, \text{m}) = (1000,\, \text{m})$. The literal
\texttt{"1000 m"\^{}\^{}}\textsf{cdt:ucum} has decimal~1000 and magnitude~1, giving
$(1000 \times 1,\, \text{m}) = (1000,\, \text{m})$. The two pairs are identical, so the
literals are value-equal.

The mapping is defined only for units that relate to their base units through a constant
multiplicative factor. The UCUM specification calls these \emph{proper units} and defines a
separate category of \emph{special units} (§21) for quantities measured on non-ratio scales,
such as temperature and logarithmic quantities. A special unit is converted to and from its
base unit by an arbitrary real function, not a constant scalar, so the \textsf{cdt:ucum}
mapping has no way to handle it. The special units identified in Section~3.3 fall into this
category: no constant factor converts Celsius to kelvin, or pH to moles per litre. An
implementation built on this mapping therefore cannot correctly evaluate equality or comparison
for special units. This is the principal known limitation of \textsf{cdt:ucum}, and it
surfaces as expected test failures in each implementation examined in this thesis.

The specification defines operator semantics for four categories. The equality operator
(\texttt{=}) tests value equality as described above. The comparison operators (\texttt{<}
and \texttt{>}) compare two \textsf{cdt:ucum} literals by their normalized real value; they are
defined only for literals of the same dimension and raise a type error when applied across
incompatible dimensions. The arithmetic operators (\texttt{+} and \texttt{-}) add or subtract
two commensurable literals, producing a result in the common base unit. The operators
(\texttt{*} and \texttt{/}) multiply or divide literals of any dimensions, combining or
cancelling dimensions accordingly, and both also accept a numeric scalar on either side. All
four categories are exposed through native SPARQL operator syntax. The specification
additionally defines the custom SPARQL function \textsf{cdt:sameDimension}, which takes two
\textsf{cdt:ucum} arguments and returns an \textsf{xsd:boolean} indicating whether they are
dimensionally commensurable.

A conforming implementation must satisfy the following obligations. It must correctly parse
all lexical forms within the \textsf{cdt:ucum} lexical space and treat forms outside it as
ill-typed. It must implement value equality using the lexical-to-value mapping. It must
implement the comparison operators with the correct dimensional restriction. It must implement
the arithmetic operators with their full type coverage, including scalar operands. It must
register \textsf{cdt:sameDimension} as a callable SPARQL function. These obligations define
the interface between the datatype and the host RDF library. Section~5 presents a methodology
for satisfying them across different library and language environments.